\documentclass[11pt]{article}
\usepackage[a4paper,margin=1in]{geometry}
\usepackage{amsmath,amssymb,amsthm}
\usepackage{booktabs}
\usepackage{graphicx}
\usepackage{xcolor}
\usepackage{hyperref}
\hypersetup{colorlinks=true,linkcolor=blue,urlcolor=blue,citecolor=blue}
\newcommand{\rt}[1]{\texttt{#1}}
\newcommand{\fn}[1]{\texttt{#1}}
\newcommand{\Tr}{\operatorname{Tr}}
\newcommand{\Phid}{\Phi^{\dagger}}

\theoremstyle{plain}
\newtheorem{proposition}{Proposition}
\newtheorem{theorem}{Theorem}
\theoremstyle{definition}
\newtheorem{definition}{Definition}
\theoremstyle{remark}
\newtheorem{remark}{Remark}

\title{\textbf{Report R22 --- what the transition graph means}\\
\large Weak components, terminal components and Hilbert-space fragmentation
in open quantum systems}
\author{Tier 1 analysis}
\date{2026-08-18}

\begin{document}
\maketitle

\begin{abstract}
Every sector count this project reports is a count of components of one
directed graph: the support graph of the channel in the computational basis.
For unitary rules the choice of component notion is immaterial, and the
project has been able to say ``Krylov sector'' without qualification. For
dissipative rules it is not immaterial, and the two natural counts can
disagree \emph{exponentially}: we have rules with a single weakly connected
component and $\psi^N$ terminal components, $\psi = 1.4656$. This report says
what the two counts are in the language of open quantum systems, and answers
the question that provoked it --- if the recurrent classes proliferate but the
weak components do not, is that Hilbert-space fragmentation? The answer is
\textbf{no}, and the reason is sharp rather than terminological. Weak
components are exactly the blocks of a \emph{strong symmetry} in the sense of
Bu\v{c}a and Prosen: the finest decomposition of the basis that makes every
Kraus operator block-diagonal, hence a genuine superselection rule with a
charge measurable at $t=0$. Terminal components are the \emph{enclosures} of
the Baumgartner--Narnhofer / Albert--Jiang asymptotic decomposition: they count
steady states, not sectors. Fragmentation is a statement about the first;
exponentially many attractors inside one block is multistability --- an
exponentially degenerate steady-state manifold with no symmetry protecting it.
The distinction is operational, not aesthetic, and we measure it: the
conserved quantities of a channel are the absorption probabilities
$\varphi_i(x)$, and a strong symmetry is precisely the statement that they are
\textbf{idempotent}. For rules whose weak components proliferate the initial
state fixes the attractor with certainty for $100\%$ of basis states; for
rules with one weak component and exponentially many attractors it does so for
$46\%$ at $N=10$, falling to $39\%$ at $N=12$ --- the rest of the time the
noise decides, not the initial condition. We map the four growth regimes onto
the literature, flag a terminology collision between two independent
``strong/weak'' dichotomies that this project uses both of, and state
precisely what a population-level graph can and cannot certify.
\end{abstract}

\tableofcontents

\section{Why the question does not arise for closed systems}
\label{sec:closed}

The object under study is the support graph $G$ of one full brick-wall cycle
$\Phi$ in the computational basis: vertices are the $2^N$ bit strings, and
$x \to y$ when the channel can take $|x\rangle\langle x|$ to something with
weight on $|y\rangle\langle y|$. Writing $\Phi(\rho) = \sum_j K_j \rho
K_j^\dagger$, this is
\begin{equation}
  x \to y \quad\Longleftrightarrow\quad
  \exists j : \langle y | K_j | x \rangle \neq 0 .
  \label{eq:graph}
\end{equation}
Three notions of ``component'' live on a directed graph: weakly connected
components (WCCs, obtained by forgetting edge directions), strongly connected
components (SCCs, mutual reachability), and terminal SCCs --- those with no
outgoing edge, which we also call recurrent classes or attractors. In general
these are different. The following says why the distinction has never had to be
made in the closed-system fragmentation literature.

\begin{proposition}[Unitarity collapses the three notions]
\label{prop:doubly}
If $\Phi(\rho) = U\rho U^\dagger$ with $U$ unitary, then every SCC of $G$ is
terminal and coincides with its WCC:
\[
  \text{WCC} \;=\; \text{SCC} \;=\; \text{terminal SCC}.
\]
\end{proposition}

\begin{proof}
The induced population dynamics is the matrix $P_{yx} = |\langle y|U|x\rangle|^2$.
Its columns sum to one because $U|x\rangle$ is normalised, and its rows sum to
one because $\langle y|U$ is; so $P$ is \emph{doubly stochastic} and the uniform
distribution is stationary with full support. In a finite Markov chain a
transient state carries zero stationary mass, so no state is transient and
every SCC is terminal. Finally, distinct terminal SCCs cannot be joined by an
edge in either direction --- an edge out of a terminal SCC would contradict
terminality --- so each is its own WCC.
\end{proof}

This is not an idle formality; it is why the project's earlier reports could
speak of ``the'' Krylov sectors. Empirically it holds exactly: across all $16$
unitary rules and all $N \in [6,16]$ the measured transient fraction is
identically $0$ and $n_{\mathrm{rec}} = n_{\mathrm{wcc}}$ at every point, and
no dissipative rule has an empty transient part at any $N$. The moment a reset
channel $D$ or $E$ enters the rule word, $P$ loses column normalisation on the
right --- two basis states are mapped onto one --- the chain acquires
transients, and the three counts separate. \textbf{The entire subject of this
report is a dissipative phenomenon with no closed-system shadow.}

\section{The dictionary}
\label{sec:dictionary}

\subsection{Weak components are a strong symmetry}

\begin{proposition}[WCCs block-diagonalise every Kraus operator]
\label{prop:wcc}
Let $\{B_\alpha\}$ be the WCCs of $G$ and $P_\alpha$ the orthogonal projector
onto $\operatorname{span}\{|x\rangle : x \in B_\alpha\}$. Then
$[K_j, P_\alpha] = 0$ for every Kraus operator $K_j$ and every $\alpha$, and
consequently $\Tr(P_\alpha \rho)$ is conserved for all time. Conversely, the
$\{B_\alpha\}$ are the \emph{finest} partition of the computational basis with
this property.
\end{proposition}

\begin{proof}
If $P_\alpha K_j P_\beta \neq 0$ for $\alpha \neq \beta$ there are
$x \in B_\beta$, $y \in B_\alpha$ with $\langle y|K_j|x\rangle \neq 0$, i.e.\ an
edge $x \to y$ in \eqref{eq:graph}; then $x$ and $y$ are weakly connected and
$\alpha = \beta$, a contradiction. Hence each $K_j = \bigoplus_\alpha P_\alpha
K_j P_\alpha$ and $P_\alpha K_j = K_j P_\alpha$. Conservation follows:
$\Tr(P_\alpha \Phi(\rho)) = \sum_j \Tr(K_j P_\alpha \rho K_j^\dagger) =
\sum_j \Tr(K_j^\dagger K_j P_\alpha \rho) = \Tr(P_\alpha\rho)$, using
$\sum_j K_j^\dagger K_j = \mathbb{1}$. Minimality holds because any coarser-than-$G$
partition with block-diagonal Kraus operators has no edges between its parts,
so it is a union of WCCs.
\end{proof}

Proposition~\ref{prop:wcc} is the whole content of the identification. In the
classification of Bu\v{c}a and Prosen~\cite{buca2012}, a symmetry of a
Lindbladian is \emph{strong} when the Hamiltonian \emph{and every jump
operator individually} commute with it, and \emph{weak} when only the
generator as a whole does. Setting $U = \sum_\alpha e^{i\theta_\alpha}
P_\alpha$ for generic phases, Proposition~\ref{prop:wcc} says exactly that the
WCC decomposition furnishes a strong symmetry, and the minimality clause says
it is the maximal one available in the computational basis. Bu\v{c}a and Prosen
establish the fact that gives this report its answer: \emph{only} a strong
symmetry implies non-trivial reductions inside the space of steady states.

So:
\begin{quote}
$n_{\mathrm{wcc}}$ counts superselection sectors. It is the open-system
successor of ``number of Krylov sectors'', and it is the count that
fragmentation is about.
\end{quote}

\subsection{Terminal components are enclosures}

The terminal SCCs are a different kind of object entirely. They are the
\emph{enclosures} of Baumgartner and Narnhofer~\cite{bn2008}, whose analysis of
GKS--Lindblad generators equips the Hilbert space with a decomposition into
mutually orthogonal subspaces and separates the transient part from the
part where probability accumulates. In the quantum-channel formulation the
same structure is the irreducible decomposition of Carbone and
Pautrat~\cite{carbone2016}: the recurrent subspace, which supports the
invariant states, decomposes into the supports of the extremal invariant
states, with the orthogonal complement being transient. Albert and
Jiang~\cite{albert2014} give the form of the asymptotic manifold,
\begin{equation}
  \rho(\infty) \;=\; \bigoplus_k \;\big( \sigma_k \otimes \rho_k \big),
  \qquad \sigma_k \in M_{n_k},
  \label{eq:aj}
\end{equation}
where $k$ runs over the enclosures, $\rho_k$ is a fixed state on the
``dissipative'' factor of enclosure $k$, and $M_{n_k}$ is a noiseless
subsystem carrying surviving coherence. Their central point --- and the one we
lean on in \S\ref{sec:answer} --- is that the infinite-time state is fixed by
the initial state together with the full list of conserved quantities.

So:
\begin{quote}
$n_{\mathrm{rec}}$ counts steady states. It is a statement about the
$t \to \infty$ limit, not about a decomposition of the Hilbert space.
\end{quote}

\subsection{The table}

\begin{center}
\small
\begin{tabular}{@{}lll@{}}
\toprule
graph object & algebraic object & physical reading \\
\midrule
WCC & block of a strong symmetry; & superselection sector; Krylov \\
    & $[K_j,P_\alpha]=0$ & sector of the dephased model \\[2pt]
$n_{\mathrm{wcc}}$ & \# blocks & \# conserved charges valid at \emph{all} times \\[2pt]
$d^{\mathrm{wcc}}_{\max}$ & largest block dimension & the R21 strong/weak-fragmentation ratio \\[4pt]
SCC & communicating class & refines the sector into transient strata \\[2pt]
transient SCCs & transient subspace & everything that decays \\[4pt]
terminal SCC & enclosure~\cite{bn2008}; irreducible & attractor; support of one extremal \\
    & component~\cite{carbone2016} & steady state \\[2pt]
$n_{\mathrm{rec}}$ & \# enclosures & \# extremal steady states \\[2pt]
$d^{\mathrm{rec}}_{\max}$ & $n_k \cdot m_k$ of the largest & noiseless-subsystem dimension \\
    & enclosure & \emph{times} mixing dimension \\
\bottomrule
\end{tabular}
\end{center}

Two structural facts follow immediately and are worth stating because they
constrain every census. First, $n_{\mathrm{rec}} \ge n_{\mathrm{wcc}}$ always,
since every weak component contains at least one terminal SCC; the inequality
is saturated exactly when each sector has a unique attractor. Second, the
familiar hyperbola $ab \ge 2$ of R9/R16 is a statement about WCCs \emph{only},
and it holds because they tile: $\sum_\alpha \dim B_\alpha = 2^N$. Recurrent
classes do not tile --- they occupy the recurrent volume $\mu^N \le 2^N$ ---
so the correct analogue is $a_{\mathrm{rec}} \, b_{\mathrm{rec}} \ge \mu$, with
$\mu = 2$ recovering the closed-system bound. For the rules discussed below
$\mu = \psi = 1.4656$ and the bound is saturated exactly.

\section{Is exponential \texorpdfstring{$n_{\mathrm{rec}}$}{n\_rec} without
exponential \texorpdfstring{$n_{\mathrm{wcc}}$}{n\_wcc} fragmentation?}
\label{sec:answer}

\subsection{The criterion}

The sharpest available definition of fragmentation is the commutant-algebra
one of Moudgalya and Motrunich~\cite{mm2022}: given the set of local terms
(Hamiltonian terms, or the local gates of a circuit), form the algebra of all
operators commuting with each of them; the system is fragmented when the
dimension of this \emph{commutant} grows exponentially in system size, as
against the polynomial growth produced by conventional symmetries such as
$U(1)$ or $SU(2)$. Their framework also supplies the classical/quantum
distinction --- ``classical'' fragmentation being fragmentation in a product
basis, which is the case here.

The open-system version replaces ``commutes with each local term'' by
``commutes with the Hamiltonian and each jump operator'', which is
Bu\v{c}a--Prosen strong symmetry, and the superoperator formulation is
developed in~\cite{mm2024}, where symmetry algebras appear as frustration-free
ground states of a local super-Hamiltonian and the strong/weak distinction is
made explicit. Under that criterion the answer is immediate:
Proposition~\ref{prop:wcc} identifies the strong-symmetry commutant in the
computational basis with the WCC decomposition, so \textbf{exponentially many
weak components is fragmentation and exponentially many terminal components is
not}. A rule with one weak component has a trivial strong-symmetry commutant:
its Hilbert space is a single Krylov sector.

\subsection{What is exponentially large instead}

This does not make the phenomenon empty --- something is exponentially large,
and it is worth naming precisely. It is the space of conserved quantities in
the Heisenberg picture. For each terminal component $T_i$ define the
\emph{absorption probability}
\begin{equation}
  \varphi_i(x) \;=\; \Pr[\,\text{absorbed in } T_i \mid \text{start in } x\,],
  \qquad
  J_i \;=\; \sum_x \varphi_i(x)\, |x\rangle\langle x| .
  \label{eq:absorb}
\end{equation}
Each $J_i$ satisfies $\Phid(J_i) = J_i$, so $\Tr(J_i \rho(t))$ is constant:
these are conserved quantities in exactly Albert and Jiang's sense, and there
are $n_{\mathrm{rec}}$ of them. Exponentially many attractors therefore do
imply exponentially many conserved quantities. The question is what \emph{kind}.

\begin{proposition}[Idempotency is the signature of a strong symmetry]
\label{prop:idem}
$J_i^2 = J_i$ for every $i$ if and only if $\varphi_i(x) \in \{0,1\}$ for
every $x$ and $i$, which holds if and only if every weak component contains
exactly one terminal component, i.e.\ $n_{\mathrm{rec}} = n_{\mathrm{wcc}}$.
In that case the $J_i$ are the $P_\alpha$ of Proposition~\ref{prop:wcc}.
\end{proposition}

\noindent This is the whole answer in one line. Fragmentation requires the
conserved quantities to be \emph{projectors} --- charges that take a definite
value on every basis state, that can be measured at $t=0$, and whose value can
never change. When $n_{\mathrm{rec}} \gg n_{\mathrm{wcc}}$ the conserved
quantities exist but are not idempotent: they are non-local, real-valued
functions on configuration space, they have no interpretation as a charge of
the initial state, and they are recovered only in the $t\to\infty$ limit.

\subsection{The operational content, measured}
\label{sec:measured}

Proposition~\ref{prop:idem} converts the abstract distinction into a number one
can compute: what fraction of basis states have their attractor determined
\emph{with certainty} by the initial condition? We built the honest population
Markov chain of the channel ($I$ keeps, $V=H$ splits $\tfrac12/\tfrac12$, $D$
and $E$ reset) and solved $\varphi = Q\varphi + R$ exactly on the transient
block.

\begin{center}
\small
\begin{tabular}{@{}llrrrr@{}}
\toprule
rule & regime & $n_{\mathrm{wcc}}$ & $n_{\mathrm{rec}}$
  & certain fate & effective \# attractors \\
\midrule
$W_{73}$  \rt{VIID}  & sectors proliferate & 41 & 41 & $100.0\%$ & $1.00$ \\
$W_{109}$ \rt{EIIV}  & sectors proliferate & 54 & 54 & $100.0\%$ & $1.00$ \\
$W_{36}$  \rt{IDDV}  & attractors only     & 10 & 60 & $72.5\%$  & $1.8$ \\
$W_{203}$ \rt{VEII}  & attractors only     & 1  & 28 & $46.1\%$  & $2.6$ \\
$W_{203}$ \rt{VEII} ($N=12$) & attractors only & 1 & 60 & $39.1\%$ & $3.2$ \\
\bottomrule
\end{tabular}
\end{center}

\noindent All rows at $N=10$ unless marked, obc0. ``Effective \# attractors''
is the participation ratio $1/\sum_i \varphi_i(x)^2$ averaged over basis
states. The dissipative rules whose \emph{weak} components proliferate behave
exactly as Proposition~\ref{prop:idem} requires: every basis state has a
certain fate, the conserved quantities are projectors, the charge is fixed at
$t=0$. $W_{203}$, with a single weak component, does not: fewer than half its
basis states have a determined fate, and the fraction \emph{falls} with $N$
while the effective number of accessible attractors rises ($\max 15.0$ at
$N=10$, $27.1$ at $N=12$).

The physical statement is therefore:

\begin{quote}
With exponentially many weak components, \textbf{the initial state decides}
which attractor the system reaches. With one weak component and exponentially
many attractors, \textbf{the noise decides}; the initial state only sets the
odds.
\end{quote}

Both are ergodicity breaking, and both give an exponentially degenerate
steady-state manifold. Only the first is Hilbert-space fragmentation.

\subsection{Where this sits in the current literature}

Two neighbouring bodies of work make the placement precise.

\emph{Dissipative freezing.} S\'anchez Mu\~noz, Buca and
co-workers~\cite{freezing2019,freezing2023} study what individual quantum
trajectories do under a \emph{strong} symmetry: initialised in a superposition
of sectors, each stochastic realisation randomly selects one and stays there,
so the conservation law that holds for the density matrix is ``broken'' at the
level of single realisations. Our case is the mirror image. There, trajectory
selection happens \emph{despite} a strong symmetry that makes the ensemble
remember; here there is no strong symmetry at all, and the trajectory-level
selection among $\psi^N$ attractors is the only structure present. The two
phenomena are easy to confuse from the trajectory picture alone and are
distinguished exactly by whether $\varphi_i$ is idempotent.

\emph{Dark states.} When the attractors are single basis states --- the
generic case among the rules where this gap appears --- the terminal
components are dark states in the standard sense: configurations immune to
every jump operator, which is precisely how the recurrent subspace is
characterised in open reaction--diffusion settings. An exponentially large
dark-state manifold reached from a structureless transient bulk is a
well-recognised target for dissipative state preparation, and it is a
different resource from a fragmented Hilbert space: one prepares states, the
other protects them.

\section{The four regimes}
\label{sec:regimes}

Writing $n_{\mathrm{wcc}} \sim a_{\mathrm{w}}^N$,
$n_{\mathrm{rec}} \sim a_{\mathrm{r}}^N$,
$d^{\mathrm{rec}}_{\max} \sim b_{\mathrm{r}}^N$, four regimes are possible and
each has a distinct reading.

\paragraph{(A/D) Sectors, attractors and attractor dimension all exponential.}
Exponentially many superselection sectors, each retaining an exponentially
large recurrent space. This is fragmentation in the full sense, and it is the
regime in which the asymptotic manifold \eqref{eq:aj} can carry an
exponentially large noiseless subsystem --- strong fragmentation is a known
route to stabilising exponentially many decoherence-free subspaces. Whether
$d^{\mathrm{rec}}_{\max} = n_k m_k$ is coherence ($n_k$) or classical mixing
($m_k$) is \emph{not} decided by the graph; see \S\ref{sec:cannot}. In our
setting the unitary rules $W_{108}$, $W_{201}$, $W_{156}$, $W_{198}$ sit here
together with the dissipative $W_{73}$, $W_{109}$.

\paragraph{(A/C) Sectors exponential, attractors trivial.} Exponentially many
sectors, each collapsing onto a small attractor --- typically one frozen
state per sector. The fragmentation is real and the memory is classical: the
long-time state is a bit string labelled by the conserved charge. This is what
dephasing does to a fragmented model, and it is the regime Li, Sala and
Pollmann~\cite{li2023} reach when they couple a fragmented system to a
dephasing bath: the quantum fragmentation is reduced to a classical one, the
stationary state becomes separable, and what survives is classical
correlations arising from fluctuations of the remaining conserved quantities.
The degenerate limiting case is $W_{204} = \rt{IIII}$, with $2^N$ sectors of
one state each.

\paragraph{(B/C) Attractors exponential, sectors not, attractors trivial.}
The regime that provoked this report. One (or polynomially many) Krylov
sectors, $\psi^N$ dark states, a transient bulk occupying $99\%+$ of the
Hilbert space, and non-idempotent conserved quantities. Not fragmentation:
\textbf{multistability with an exponentially degenerate dark-state manifold}.
The exemplar is $W_{203} = \rt{VEII}$, whose entire $2^N$-dimensional space is
one weak component.

\paragraph{(B/D) Attractors exponential and exponentially large, sectors not.}
Exponentially many enclosures each exponentially large, inside a single
superselection sector. This would be the most interesting regime --- an
exponentially degenerate manifold of exponentially large noiseless
subsystems with no symmetry organising them --- and it is \emph{empty} in the
256-rule family at obc0. We record its absence as an observation, not a
theorem; nothing we know forbids it, and the volume constraint
$a_{\mathrm{r}} b_{\mathrm{r}} \le 2$ leaves room for it.

\section{A terminology collision worth flagging}
\label{sec:collision}

This project now uses two independent dichotomies that are both called
``strong versus weak'', and they are orthogonal. Conflating them would
invalidate R21's conclusions or this report's, depending on the direction of
the error.

\begin{center}
\small
\begin{tabular}{@{}p{0.45\textwidth}p{0.45\textwidth}@{}}
\toprule
\textbf{strong / weak \emph{fragmentation}} &
\textbf{strong / weak \emph{symmetry}} \\
Sala \emph{et al.}~\cite{sala2020}, Khemani \emph{et al.}~\cite{khemani2020} &
Bu\v{c}a--Prosen~\cite{buca2012} \\
\midrule
About the \emph{size} of the largest sector relative to the symmetry sector
containing it: strong iff $D_{\max}(s)/\dim s \to 0$. &
About \emph{how} a symmetry acts: strong iff every jump operator commutes with
it individually, weak iff only the generator does. \\[4pt]
A question you ask \emph{after} you have a sector decomposition. &
A question about whether you have a sector decomposition at all. \\[4pt]
Used in R21, which grades 108/201/156/198 as strongly fragmented. &
Used here, to say what $n_{\mathrm{wcc}}$ counts. \\
\bottomrule
\end{tabular}
\end{center}

A model can be weakly fragmented with a perfectly good strong symmetry
($W_{60}$, $W_{102}$ in R21), and --- as \S\ref{sec:answer} shows --- can have
exponentially many steady states with no strong symmetry whatever. There is
also a third, newer use of the same words in the mixed-state literature:
\emph{strong-to-weak spontaneous symmetry breaking}~\cite{swssb2024}, where a
strong symmetry of a density matrix degrades to a weak one and the order
parameter is a fidelity correlator rather than a two-point function. That is a
statement about phases of mixed states, not about sector counting, and it
should not be read into either column above. Whether the boundary between our
regimes A and B can be crossed as a genuine SWSSB transition in a family of
channels is a natural question and one we have not touched.

\section{What a population graph cannot certify}
\label{sec:cannot}

The support graph \eqref{eq:graph} is the \emph{population} skeleton of the
channel. It records which basis states feed which, and nothing about the
phases. Three consequences bound every claim in this report and in the
project's dissipative reports generally.

\begin{enumerate}
\item \textbf{It is a lower bound on fragmentation.} A channel may have
conserved quantities invisible in the computational basis --- the fragmentation
may be quantum, i.e.\ in a non-product basis, in the sense
of~\cite{mm2022}. The WCC count is then an undercount. This is not
hypothetical for us: R-T15 identified the single-$V$ commutant of
$W_{108}/W_{201}/W_{156}/W_{198}$ as a Temperley--Lieb algebra, and the
Temperley--Lieb model is exactly the vehicle Li, Sala and Pollmann use for
\emph{quantum} HSF in an entangled basis~\cite{li2023}. Their
fragmentation-preserving coupling produces a steady state with coherent memory
carried by an extensive set of \emph{non-commuting} conserved quantities. No
diagonal graph can see any of that.

\item \textbf{It cannot split $d^{\mathrm{rec}}_{\max}$ into $n_k \cdot m_k$.}
The graph gives the dimension of an enclosure's support; \eqref{eq:aj} says
that dimension factorises into a noiseless subsystem and a mixing factor, and
which one carries the growth is a question about coherence. The project has
answered it in both directions already: R17 found the exponential terminal
SCCs to be dense, mixing digraphs rather than cycles --- so $m_k$ is large and
$n_k$ small, a classical attractor --- while the $X$-gate correction
established the opposite for permutation circuits, where the functional graph
describes only populations and exponentially large decoherence-free subspaces
sit underneath it. \emph{The graph alone never settles this.}

\item \textbf{The equality $\dim(\text{fixed points of } \Phid) =
n_{\mathrm{rec}}$ is a classical statement.} For the population chain the
conserved quantities are exhausted by the absorption probabilities
\eqref{eq:absorb}. The full channel can carry additional fixed points built
from coherences between enclosures of equal structure, which is precisely the
$M_{n_k}$ factor in \eqref{eq:aj}. So $n_{\mathrm{rec}}$ is a lower bound on
the steady-state count too.
\end{enumerate}

\section{Summary}

\begin{itemize}
\item WCCs are the blocks of a \emph{strong symmetry}: the finest basis
decomposition making every Kraus operator block-diagonal
(Prop.~\ref{prop:wcc}). Their count is the open-system Krylov sector count,
and exponential growth of it is Hilbert-space fragmentation.
\item Terminal SCCs are \emph{enclosures}: supports of extremal steady states.
Their count is a $t\to\infty$ statement and is always $\ge n_{\mathrm{wcc}}$.
\item Exponentially many terminal SCCs inside one WCC is \textbf{not}
fragmentation. It is multistability: the conserved quantities are absorption
probabilities, which fail to be idempotent (Prop.~\ref{prop:idem}), and the
attractor is selected by noise rather than fixed by the initial state --- a
$100\%$ versus $46\%$ effect in direct measurement (\S\ref{sec:measured}).
\item The distinction has no closed-system analogue: unitarity makes the
population dynamics doubly stochastic, which kills the transient part and
collapses all three component notions (Prop.~\ref{prop:doubly}).
\item ``Strong/weak'' means two different things in this project, and now a
third thing in the mixed-state literature (\S\ref{sec:collision}).
\item Everything here is certified at the level of populations only
(\S\ref{sec:cannot}); coherent fragmentation would be invisible, and our own
Temperley--Lieb structure is a live candidate for it.
\end{itemize}

\section*{Reproduce}
\addcontentsline{toc}{section}{Reproduce}

The measurement of \S\ref{sec:measured} needs transition \emph{probabilities},
which the project's component engine does not carry: \fn{graph/scc.py} works
with the unweighted successor relation, since that is all that component
structure requires. The weighted one-cycle chain is therefore rebuilt from the
gate semantics ($I$ keeps, $V=H$ splits $\tfrac12/\tfrac12$, $D$ and $E$
reset), in \fn{scaling/absorption.py}, at cost $O(4^N)$. It is a small-$N$
diagnostic, not a sweep.

\begin{verbatim}
    python -m qca_fragmentation.scaling.absorption --table
    pytest qca_fragmentation/tests/test_absorption.py     # 58 tests
\end{verbatim}

\noindent The tests pin every number in the table of \S\ref{sec:measured}, and
check both directions of Proposition~\ref{prop:idem} as well as the double
stochasticity of Proposition~\ref{prop:doubly} --- the latter holding for the
unitary rules and failing for the dissipative ones.

\begin{thebibliography}{99}

\bibitem{buca2012} B.~Bu\v{c}a and T.~Prosen,
\emph{A note on symmetry reductions of the Lindblad equation: transport in
constrained open spin chains},
New J.\ Phys.\ \textbf{14}, 073007 (2012);
\href{https://arxiv.org/abs/1203.0943}{arXiv:1203.0943}.

\bibitem{bn2008} B.~Baumgartner and H.~Narnhofer,
\emph{Analysis of quantum semigroups with GKS--Lindblad generators: II.
General},
J.\ Phys.\ A \textbf{41}, 395303 (2008);
\href{https://arxiv.org/abs/0806.3164}{arXiv:0806.3164}.

\bibitem{albert2014} V.~V.~Albert and L.~Jiang,
\emph{Symmetries and conserved quantities in Lindblad master equations},
Phys.\ Rev.\ A \textbf{89}, 022118 (2014);
\href{https://arxiv.org/abs/1310.1523}{arXiv:1310.1523}.

\bibitem{carbone2016} R.~Carbone and Y.~Pautrat,
\emph{Irreducible decompositions and stationary states of quantum channels},
Rep.\ Math.\ Phys.\ \textbf{77}, 293 (2016);
\href{https://arxiv.org/abs/1507.08404}{arXiv:1507.08404}.

\bibitem{mm2022} S.~Moudgalya and O.~I.~Motrunich,
\emph{Hilbert space fragmentation and commutant algebras},
Phys.\ Rev.\ X \textbf{12}, 011050 (2022);
\href{https://arxiv.org/abs/2108.10324}{arXiv:2108.10324}.

\bibitem{mm2024} S.~Moudgalya and O.~I.~Motrunich,
\emph{Symmetries as ground states of local superoperators: hydrodynamic
implications},
PRX Quantum \textbf{5}, 040330 (2024);
\href{https://arxiv.org/abs/2309.15167}{arXiv:2309.15167}.

\bibitem{li2023} Y.~Li, P.~Sala and F.~Pollmann,
\emph{Hilbert space fragmentation in open quantum systems},
Phys.\ Rev.\ Research \textbf{5}, 043239 (2023);
\href{https://arxiv.org/abs/2305.06918}{arXiv:2305.06918}.

\bibitem{sala2020} P.~Sala, T.~Rakovszky, R.~Verresen, M.~Knap and F.~Pollmann,
\emph{Ergodicity breaking arising from Hilbert space fragmentation in
dipole-conserving Hamiltonians},
Phys.\ Rev.\ X \textbf{10}, 011047 (2020).

\bibitem{khemani2020} V.~Khemani, M.~Hermele and R.~Nandkishore,
\emph{Localization from Hilbert space shattering: from theory to physical
realizations},
Phys.\ Rev.\ B \textbf{101}, 174204 (2020).

\bibitem{freezing2019} C.~S\'anchez Mu\~noz, B.~Buca, J.~Tindall,
A.~Gonz\'alez-Tudela, D.~Jaksch and D.~Porras,
\emph{Symmetries and conservation laws in quantum trajectories: dissipative
freezing},
Phys.\ Rev.\ A \textbf{100}, 042113 (2019).

\bibitem{freezing2023} J.~Tindall, D.~Jaksch and C.~S\'anchez Mu\~noz,
\emph{On the generality of symmetry breaking and dissipative freezing in
quantum trajectories},
SciPost Phys.\ Core \textbf{6}, 004 (2023);
\href{https://arxiv.org/abs/2204.06585}{arXiv:2204.06585}.

\bibitem{swssb2024} L.~A.~Lessa, R.~Ma, J.-H.~Zhang, Z.~Bi, M.~Cheng and
C.~Wang,
\emph{Strong-to-weak spontaneous symmetry breaking in mixed quantum states},
PRX Quantum \textbf{6}, 010344 (2025);
\href{https://arxiv.org/abs/2405.03639}{arXiv:2405.03639}.

\end{thebibliography}

\end{document}
